\documentclass[10pt,a4paper]{amsart}
\usepackage[margin=19mm]{geometry}
\usepackage{amsmath,amssymb,mathtools}
\usepackage[scr=rsfs,cal=euler]{mathalfa}
\usepackage{xfrac}
\usepackage[unicode,hidelinks,bookmarks=false]{hyperref}
\newcommand{\ie}{i.e.\ }
\newcommand{\cf}{cf.\ }
\newcommand{\resp}{respectively\ }
\newcommand{\Subsection}{\subsection}
\providecommand{\nobreakdash}{\nobreak\hbox{-}}
\setlength{\emergencystretch}{2em}
\multlinegap0pt
\usepackage{mathtools}
\usepackage{amssymb}
\usepackage{xcolor}
\usepackage[shortlabels]{enumitem}
\newtheorem{prop}{Proposition}
\newcommand{\Div}{{\rm Div}}
\newcommand{\NS}{{\operatorname{NS}}}
\newcommand{\Q}{{\mathbb Q}}
\newcommand{\Z}{{\mathbb Z}}
\newcommand{\pr}{{\rm{pr}}}
\title{Zero-cycles on quasi-projective surfaces over $p$-adic fields}
\author{Evangelia Gazaki, Jitendra Rathore}
\date{Source: arXiv:2507.20076v2 (CC BY 4.0)}
\begin{document}
\maketitle

\noindent\emph{Source and licence.} Zero-cycles on quasi-projective surfaces over $p$-adic fields. Version arXiv:2507.20076v2 (CC BY 4.0), distributed by arXiv under the Creative Commons Attribution 4.0 International licence, http://creativecommons.org/licenses/by/4.0/, as recorded in the arXiv licence metadata for this exact version and shown on the abstract page https://arxiv.org/abs/2507.20076v2. Full licence notice: \texttt{licenses/arxiv-2507.20076-CC-BY-4.0.txt}.\par\smallskip

\noindent\textit{Editorial scope.} Counted unit: the article's prop prop:NSXcomputations (source0.tex lines 913--914). The article's complete original proof of it is delivered with it (source0.tex lines 915--959, one \texttt{proof} environment of 45 lines). No mathematics is written, repaired or re-derived: the statement and the proof are the article's own lines, in the article's own notation, and no retained line is substituted or re-flowed.  No further source-local context is needed: every cross-reference of every retained line resolves inside the two delivered spans.  Nothing was omitted from any retained span, so the package carries no omission record.  No retained line contains a citation, so the wrapper carries no bibliography and no bibliography entry.  No retained line contains a figure environment, an image-inclusion command or drawing code, so no image is delivered and the package declares no asset dependency.  Editorial formatting only, in addition: an AMS \texttt{amsart} preamble that restates the article's own declarations and macros used by the retained lines, namely the environments \texttt{prop} on separate counters, together with the standard packages those lines need.  The retained environments print in this document as Proposition 1 in order of appearance, whereas the article prints the same items with the numbering of its own full document.  The complete native source of the article as distributed in the arXiv e-print for this exact version is bundled unchanged as \texttt{sources/182260/source0.tex}.
\begin{prop}\label{prop:NSXcomputations} Consider the set-up of the previous paragraph. Then $\iota^\star(\NS(X))$ is a subgroup of $\Z^r$ of rank $n$. 
\end{prop}

\begin{proof} 
By assumption the elements $[D_1],\ldots, [D_n]$ are $\Z-$linearly independent in $\NS(X)$. For each $m=1,\ldots, n$ we denote by $\pr_m:\Z^r\to\Z^m$ the projection onto the first $m$ factors. We consider the composition 
\[f_m:\NS(X)\xrightarrow{\iota^\star}\Z^r\xrightarrow{\pr_m}\Z^m.\]
We claim that this map has finite cokernel for every $1\leq m\leq n$. 
We proceed by induction on $m\geq 1$. Since $[D_1]\in \NS(X)$ is $\Z$-linearly independent, it is not numerically equivalent to $0$, and hence there is some curve $C$ in $X$ such that $C\cdot D_1\neq 0$. This proves the $m=1$ case. 

Next suppose $2\leq m\leq n$. By induction hypothesis there are divisors $S_1,\ldots, S_{m-1}$ in $X$ such that $\{\pr_{m-1}\circ\iota^\star([S_i]):1\leq i\leq m-1\}$ is a $\Z$-linearly independent subset of $\Z^{m-1}$. For each $i, j\in\{1,\ldots, m-1\}$ set $\mu_{ij}:=S_i\cdot D_j$. It follows that the matrix $A=(\mu_{ij})\in M_{m-1}(\Q)$ is invertible. It is clear that the set $\{f_{m}([S_1]),\ldots, f_{m}([S_{m-1}])\}\subset\Z^{m}$ is also $\Z$-linearly independent as the first $m-1$ coordinates of these vectors are the same as before. This means that $f_m(\NS(X))$ is a subgroup of $\Z^m$ of rank at least $m-1$. Suppose for contradiction that $\pr_m\circ\iota^\star$ does not have finite cokernel. This means that $\{f_{m}([S_1]),\ldots, f_{m}([S_{m-1}])\}\subset\Z^{m}$ is a maximal $\Z$-linearly independent subset of $f_m(\NS(X))$. It then follows that for every curve $C$ in $X$ there exist $q_{1,C},\ldots, q_{m-1,C}\in\Q$ such that 
\begin{equation}\label{eq10}
    (C\cdot D_1,\ldots, C\cdot D_m)=\sum_{i=1}^{m-1} q_{i,C}(S_i\cdot D_1,\ldots, S_i\cdot D_m).
\end{equation}
Let us rewrite the first $m-1$ coordinates of (\ref{eq10}) as follows: 
\[
\left\{\begin{array}{c}
    C\cdot D_1=\sum_{i=1}^{m-1} q_{i,C}\mu_{i1}   \\
     \cdots\\
     C\cdot D_{m-1}=\sum_{i=1}^{m-1} q_{i,C}\mu_{i,m-1}
\end{array}\right.
\] This can be thought of as a $\Q$-linear system with $m-1$ equations in the unknowns $q_{1,C},\ldots, q_{m-1,C}$. The matrix of this system is the transpose of $A$. Since this matrix is invertible, it follows that the system has a unique solution. In particular, for each $i=1,\ldots, m-1$, $q_{i,C}$ has an expression 
\begin{equation}
    q_{i,C}=\sum_{j=1}^{m-1}\lambda_{ij} (C\cdot D_j),
\end{equation} for some $\lambda_{ij}\in\Q$, which are independent of $C$. 

For $i\in\{1,\ldots, m-1\}$ set $\gamma_i=S_i\cdot D_m$, which is again independent of $C$. 
Combining all the above, the last coordinate of (\ref{eq10}) can be rewritten as 
\begin{eqnarray*}
 C\cdot D_m=&&\sum_{i=1}^{m-1} q_{i,C}(S_i\cdot D_m)=\sum_{i=1}^{m-1} q_{i,C}\gamma_i\\
 =&&\sum_{i=1}^{m-1}\gamma_i\left(\sum_{j=1}^{m-1}\lambda_{ij} (C\cdot D_j)\right).
\end{eqnarray*}
If we expand this relation we potentially reach an expression of the form
\[C\cdot\left(D_m-\sum_{i=1}^{m-1}\alpha_i D_i\right)=0,\] for some $\alpha_i\in\Q$ independent of $C$. Write $\alpha_i=\frac{a_i}{b}$ for some $a_i, b\in\Z$ with $b\neq 0$. Since the curve $C$ was arbitrary, it follows that the divisor $\displaystyle bD_m-\sum_{i=1}^{m-1}a_i D_i$ is numerically equivalent to $0$. But this contradicts the $\Z$-linear independence of $[D_1],\ldots, [D_m]$ for $m\leq n$ in $\NS(X)$. We conclude that $f_m(\NS(X))$ is a subgroup of $\Z^m$ of rank $m$, which completes the induction.




To finish the proof we need to show that the rank of $\iota^\star(\NS(X))$ is exactly $n$. For simplicity of notation we will only prove that $f_{n+1}(\NS(X))\otimes\Q$ is an $n$-dimensional subspace of $\Q^{r}$. The argument for a general $n+1\leq m\leq r$ is similar. By assumption, the subset $\{[D_1],\ldots, [D_n]\}\subset\NS(X)$ is $\Z$-linearly independent, while $\{[D_1],\ldots, [D_n], [D_{n+1}]\}$ is not. Working in $\NS(X)\otimes\Q$, this means that we can find $\lambda_1,\ldots, \lambda_n\in\Q$ such that $\displaystyle D_{n+1}-\sum_{i=1}^n \lambda_i D_i$ is numerically equivalent to $0$. Let $S\in\Div(X)$.  Then we can write
\begin{eqnarray*}
    f_{n+1}(S)=&&(S\cdot D_1,\ldots, S\cdot D_n, S\cdot D_{n+1})=\left(S\cdot D_1,\ldots, S\cdot D_n, \sum_{i=1}^n \lambda_i (S\cdot D_{i})\right)\\
    =&& S\cdot D_1(1,0,\ldots, 0, \lambda_1)+\cdots + S\cdot D_n (0,\ldots, 1, \lambda_n)\\
    &&=c_1(1,0,\ldots, 0, \lambda_1)+\cdots+c_n(0,\ldots, 1, \lambda_n),
\end{eqnarray*} where we set $c_i=S\cdot D_i$ for $i=1,\ldots, n$. Since the $\lambda_i\in\Q$ are independent of $S$, it follows that the vectors $\{(1,0,\ldots, \lambda_1),\ldots (0,\ldots, 1, \lambda_n)\}$ span $f_{n+1}(\NS(X))$, from where the claim follows. 




\end{proof}

\end{document}
